\chapter{Ad\`eles, Id\`eles and Tate's Thesis}
\begin{center}\itshape
Riemann proved the functional equation of $\zeta$ with the Poisson summation formula.\\
Tate proved it with the Poisson summation formula too,\\
but at every prime at once.
\end{center}

In 1859 Riemann wrote his only paper on number theory, eight pages long. Its second proof of the functional equation
$\Lambda(s)=\Lambda(1-s)$, the one we followed in Tome~4, rests on a single identity, $\theta(1/y)=\sqrt y\,\theta(y)$, which is
the Poisson summation formula applied to the Gaussian. Sixty years later Hecke extended the proof to the zeta function of every
number field and to a zoo of $L$-functions attached to ``Gr\"o\ss encharaktere''. Hecke's proofs work, but they are long and
technical: theta functions in several variables, units, ideal classes, each handled separately, each a small miracle.

Then the point of view changed. In 1936 Claude Chevalley, working on class field theory, decided that the ideals were getting in the
way and replaced them by \emph{id\`eles} (``ideal elements''): sequences $(x_v)$ indexed by all the places of the field, one
component in each completion. Andr\'e Weil added the additive version, the \emph{ad\`eles}. And in 1950 a Princeton graduate student,
John Tate, wrote a thesis under Emil Artin with an outrageous idea: do Fourier analysis \emph{on the ad\`eles}. The ring $\mathbb{A}$ of
ad\`eles of $\Q$ contains $\Q$ as a discrete subgroup, as $\R$ contains $\Z$, and the quotient $\mathbb{A}/\Q$ is compact, as $\R/\Z$ is.
So there is a Poisson summation formula $\sum_{q\in\Q}f(q)=\sum_{q\in\Q}\hat f(q)$, and the functional equation of every Hecke
$L$-function falls out of it in a page. The Gamma factor that looked like a decoration in Riemann's formula turns out to be the
Euler factor at the prime $\infty$, the same kind of object as $(1-p^{-s})^{-1}$. (Kenkichi Iwasawa found the same proof
independently at about the same time; it became known through a 1952 letter to Dieudonn\'e, published only in 1992.)

This chapter does Tate's thesis for $\Q$ in full, with all measure theory built by hand, and proves the functional equations of
$\zeta$ and of every Dirichlet $L$-function. Then it turns to a number field $K$ and proves Fujisaki's theorem, that the norm-one
id\`eles modulo $K^\times$ form a compact group, from Minkowski's theorem; the finiteness of the class number and Dirichlet's unit
theorem drop out as two shadows of that single compactness. Tome~1 promised, at the end of its chapter on $p$-adic numbers, that
the modern point of view treats all completions of $\Q$ at once. This is that point of view.

\begin{reminder}[from Tomes 1, 4, 5, 6 and Chapters 3--6]
\emph{$p$-adic numbers} (Tome~1, Chapter~12): $|x|_p=p^{-v_p(x)}$, $\Z_p=\{|x|_p\le1\}$ is compact and $\Z$ is dense in it
(Tome~1, Proposition~12.4); two balls of $\Q_p$ are disjoint or nested. \emph{Fourier analysis on $\R$} (Tome~4, Chapter~8):
$\hat f(\xi)=\int f(x)e^{-2\pi ix\xi}\dd x$; $e^{-\pi x^2}$ is its own transform (Tome~4, Theorem~8.3); if $xf(x)$ is integrable,
$\hat f\,'$ is the transform of $-2\pi ixf(x)$ (Tome~4, Proposition~8.2); inversion $\hat{\hat f}(x)=f(-x)$ (Tome~4, Theorem~8.6);
Poisson summation $\sum_nf(n)=\sum_m\hat f(m)$ when $f$ and $\hat f$ are continuous and $O((1+|x|)^{-2})$ (Tome~4, Theorem~8.9).
\emph{Riemann's functional equation} (Tome~4, Theorem~10.8) and holomorphic parameter integrals (Tome~4, Theorem~9.15).
$\Gamma(s)=\int_0^\infty t^{s-1}e^{-t}\dd t$, and the duplication formula (Tome~4, Theorem~5.14). \emph{Gauss sums} (Tome~5,
Theorem~3.5): for $\chi$ primitive modulo $N$ and every integer $m$, $\sum_{c=1}^N\chi(c)e^{2\pi icm/N}=\overline{\chi(m)}\tau(\chi)$.
\emph{Measure theory} (Tome~6): Carath\'eodory's extension theorem from a premeasure on a semiring, with uniqueness (Tome~6,
Theorem~6.10); Tychonoff's theorem (Tome~6, Theorem~4.11). From this book: Minkowski's theorem (Theorem~\ref{thm:minkowski}),
$\operatorname{covol}j(I)=2^{-r_2}\sqrt{|d_K|}N(I)$ (Proposition~\ref{prop:covol}), the log embedding (Chapter~4), the places of $K$
and the product formula $\prod_v|x|_v=1$ (Proposition~\ref{prop:productformula}).
\end{reminder}

\section{The ad\`eles of \texorpdfstring{$\Q$}{Q}}

The \emph{places} of $\Q$ are the primes $p$ and the symbol $\infty$, with $\Q_\infty=\R$ and $|x|_\infty=|x|$. A rational number is
an integer at all but finitely many primes. This suggests the following definition.

\begin{definition}
The ring of \emph{finite ad\`eles} is
\[
\mathbb{A}_f=\Bigl\{(x_p)_p\in\prod_p\Q_p:\ x_p\in\Z_p\text{ for all but finitely many }p\Bigr\},
\]
with componentwise operations. The ring of \emph{ad\`eles} is $\mathbb{A}=\R\times\mathbb{A}_f$. We write $\hat{\Z}=\prod_p\Z_p\subseteq\mathbb{A}_f$.
\end{definition}

Sums and products of two ad\`eles are ad\`eles, since the union of two finite sets of exceptional primes is finite. We put on $\mathbb{A}_f$
the topology generated by the sets
\[
\prod_{p\in S}U_p\times\prod_{p\notin S}\Z_p,\qquad S\text{ finite},\ U_p\subseteq\Q_p\text{ open},
\]
(they are stable under finite intersections, so they form a basis), and on $\mathbb{A}$ the product topology with $\R$. This is
\emph{not} the topology induced by the product $\prod_p\Q_p$: it is finer, and $\hat{\Z}$ is \emph{open} in $\mathbb{A}_f$. On $\hat{\Z}$ itself the
two topologies agree, so $\hat{\Z}$ is compact by Tychonoff's theorem; hence $\mathbb{A}_f$ and $\mathbb{A}$ are locally compact. Addition and
multiplication are continuous, because they are on each $\Q_p$ and they preserve $\prod_{p\notin S}\Z_p$. For a positive rational
$N$ the set $N\hat{\Z}$ is an open compact subgroup, and these subgroups form a basis of neighbourhoods of $0$: every basic open set
containing $0$ contains $\prod_{p\in S}p^k\Z_p\times\prod_{p\notin S}\Z_p=N\hat{\Z}$ for $N=\prod_{p\in S}p^k$.

The \emph{diagonal embedding} $\Q\to\mathbb{A}$, $q\mapsto(q,q,q,\dots)$, makes sense since $q\in\Z_p$ for every $p$ not dividing its
denominator. It is a ring homomorphism, and we write $\Q$ for its image. By the Chinese remainder theorem $\Z$ is dense in $\hat{\Z}$:
given $z\in\hat{\Z}$ and $N=\prod p^{k_p}$, there is $n\in\Z$ with $n\equiv z_p\pmod{p^{k_p}}$ for all $p\mid N$, that is
$z-n\in N\hat{\Z}$.

For $x\in\Q_p$ let $\{x\}_p$ be its \emph{fractional part}: the unique element of $\Z[\frac1p]\cap\left[0,1\right[$ with
$x-\{x\}_p\in\Z_p$. (Write $x=\sum_{k\ge-m}a_kp^k$ with digits $0\le a_k<p$; then $\{x\}_p=\sum_{k<0}a_kp^k$. Uniqueness: two
candidates differ by an element of $\Z_p\cap\Z[\frac1p]=\Z$ lying in $\left]-1,1\right[$.)

\begin{lemma}\label{lem:AfQZ}
$\mathbb{A}_f=\Q+\hat{\Z}$, and $\Q\cap\hat{\Z}=\Z$. More precisely, for $x\in\mathbb{A}_f$ the rational number $r=\sum_p\{x_p\}_p$ (a finite sum)
satisfies $x-r\in\hat{\Z}$.
\end{lemma}
\begin{proof}
The sum is finite since $\{x_p\}_p=0$ when $x_p\in\Z_p$. Fix a prime $q$. Then $x_q-r=(x_q-\{x_q\}_q)-\sum_{p\ne q}\{x_p\}_p$, and
each term lies in $\Z_q$: the first by definition, the others because $\Z[\frac1p]\subseteq\Z_q$ for $p\ne q$. So $x-r\in\hat{\Z}$. A
rational number lying in $\Z_p$ for every $p$ has no prime in its denominator, so $\Q\cap\hat{\Z}=\Z$.
\end{proof}

This is nothing but the decomposition of a rational number into partial fractions: $\frac{5}{12}=\frac34+\frac23-1$, with
$\frac34=\{\frac5{12}\}_2$ and $\frac23=\{\frac5{12}\}_3$.

\begin{theorem}\label{thm:Qdiscrete}
\begin{enumerate}[label=(\alph*)]
\item $\Q$ is a discrete subgroup of $\mathbb{A}$.
\item Every ad\`ele can be written uniquely as $q+d$ with $q\in\Q$ and $d\in D=\left[0,1\right[\times\hat{\Z}$.
\item The map $\R\times\hat{\Z}\to\mathbb{A}/\Q$ is surjective with kernel the diagonal copy of $\Z$, and $\mathbb{A}/\Q$ is compact.
\end{enumerate}
\end{theorem}
\begin{proof}
(a) The open set $\left]-1,1\right[\times\hat{\Z}$ meets $\Q$ in the integers $n$ with $|n|<1$ (Lemma~\ref{lem:AfQZ}), that is in $\{0\}$.
Translating, every $q\in\Q$ has a neighbourhood containing no other rational.

(b) Let $x=(x_\infty,x_f)$. By Lemma~\ref{lem:AfQZ}, $x_f=r+z$ with $r\in\Q$, $z\in\hat{\Z}$. Let $n=\lfloor x_\infty-r\rfloor$; then
$x-r-n=(x_\infty-r-n,\ z-n)\in D$. If $q+d=q'+d'$, then $q-q'=d'-d\in\Q\cap\bigl(\left]-1,1\right[\times\hat{\Z}\bigr)=\{0\}$.

(c) Surjectivity is (b). An element $(y,z)\in\R\times\hat{\Z}$ lies in $\Q$ iff it is a rational $q$ with $q\in\hat{\Z}$, that is $q\in\Z$.
Finally $\mathbb{A}/\Q$ is the image of the compact set $[0,1]\times\hat{\Z}$ under the continuous projection, so it is compact.
\end{proof}

So $\mathbb{A}/\Q\cong(\R\times\hat{\Z})/\Z$, a compact connected group called the \emph{solenoid}. Think of it as a circle $\R/\Z$ that
has been unrolled into a profinite family of strands: going once around the circle moves you by $-1$ in $\hat{\Z}$, so you never close up.
Figure~\ref{fig:solenoid} draws the $2$-adic version.

\begin{figure}[ht]
\centering
\fig{ch09-solenoid}
\caption{The $2$-adic solenoid $(\R\times\Z_2)/\Z$, a picture of $\mathbb{A}/\Q$ with one prime. The tree on the left is $\Z_2$ (branch
$0$ or $1$ at each digit); each strand is $\left[0,1\right[\times(a+8\Z_2)$, and the end $(1,a)$ of a strand is glued to the start
$(0,a-1)$ of the next. Going around once subtracts $1$ in $\Z_2$, which changes infinitely many digits only when $a=0$.}
\label{fig:solenoid}
\end{figure}

\begin{think}
$\Q$ is dense in $\R$, and dense in every $\Q_p$. How can it be \emph{discrete} in $\mathbb{A}$? A rational close to $0$ in $\mathbb{A}$ must be
small at $\infty$ \emph{and} integral at every prime simultaneously, and the only such rational is $0$. Discreteness is a
global statement; it is the product formula in disguise (a non-zero rational cannot be small everywhere). Exercise~\ref{ex:strongapprox}
shows the opposite phenomenon: drop the place $\infty$ and $\Q$ becomes dense in $\mathbb{A}_f$.
\end{think}

\section{Haar measures}

To do Fourier analysis we need to integrate. On $\R$ we have Lebesgue measure; on $\Q_p$ we build its analogue.

\begin{proposition}[Haar measure on $\Q_p$]\label{prop:haarQp}
There is a unique measure $\mu_p$ on the Borel sets of $\Q_p$ with $\mu_p(a+p^k\Z_p)=p^{-k}$ for all $a\in\Q_p$ and $k\in\Z$.
It is translation invariant, and $\mu_p(aE)=|a|_p\,\mu_p(E)$ for $a\in\Q_p^\times$. Consequently
$\int g(ax)\dd\mu_p(x)=|a|_p^{-1}\int g\dd\mu_p$ for every integrable or non-negative measurable $g$.
\end{proposition}
\begin{proof}
Let $\mathcal S$ be the family made of $\varnothing$ and the balls $a+p^k\Z_p$. It is a semiring. Two balls are disjoint or
nested, so the intersection of two members is a member. If $B'\subseteq B$ are balls of radii $p^{-k'}\le p^{-k}$, then $B$ is the
disjoint union of the $p^{k'-k}$ cosets of $p^{k'}\Z_p$ that it contains, one of which is $B'$, so $B\setminus B'$ is a finite disjoint
union of balls; and if $B'\not\subseteq B$, then $B\setminus B'$ is $B$ or $\varnothing$.

Put $\mu_0(a+p^k\Z_p)=p^{-k}$, $\mu_0(\varnothing)=0$. \emph{Finite additivity.} Suppose a ball $B$ of radius $p^{-k}$ is the disjoint
union of balls $B_1,\dots,B_m$ of radii $p^{-k_i}$. Let $K=\max k_i$. Each $B_i$ is the disjoint union of $p^{K-k_i}$ balls of radius
$p^{-K}$, each of mass $p^{-K}$, so $\mu_0(B_i)$ is $p^{-K}$ times the number of these small balls; likewise for $B$, which is made of
the same small balls, each counted once. So $\mu_0(B)=\sum\mu_0(B_i)$. \emph{Countable additivity.} If $B=\bigsqcup_{i\ge1}B_i$ with
non-empty $B_i\in\mathcal S$, then, as $B$ is compact (homeomorphic to $\Z_p$) and the $B_i$ are open, finitely many of them, say
$B_1,\dots,B_m$, cover $B$; any other $B_i$ would be disjoint from $B_1\cup\dots\cup B_m=B$, hence empty. So the sum is finite and
finite additivity applies.

By Carath\'eodory's theorem (Tome~6, Theorem~6.10), $\mu_0$ extends to a measure on $\sigma(\mathcal S)$, uniquely since
$\Q_p=\bigcup_kp^{-k}\Z_p$ is an increasing union of members of finite mass. Every open set is a union of balls with centres in the
countable set $\Z[\frac1p]$ (dense in $\Q_p$) and radii $p^{-k}$, hence a countable union of members of $\mathcal S$; so $\sigma(\mathcal S)$
is the Borel $\sigma$-algebra.

For $b\in\Q_p$, $E\mapsto\mu_p(b+E)$ is a measure that agrees with $\mu_p$ on $\mathcal S$, so it equals $\mu_p$ by uniqueness. For $a\ne0$,
$E\mapsto|a|_p^{-1}\mu_p(aE)$ is a measure, and $a(b+p^k\Z_p)=ab+p^{k+v_p(a)}\Z_p$ has mass $p^{-k}|a|_p$, so it also agrees with
$\mu_p$ on $\mathcal S$. The integral formula holds for $g=\mathbf 1_E$, since $\mathbf 1_E(ax)=\mathbf 1_{a^{-1}E}(x)$; hence for simple
functions, then for non-negative ones by monotone convergence, then for integrable ones.
\end{proof}

We write $\dd x$ for $\mu_p$. A first computation: $\Z_p\setminus\{0\}$ is the disjoint union of the shells $p^k\Z_p^\times$, $k\ge0$,
of mass $p^{-k}-p^{-k-1}=p^{-k}(1-\frac1p)$, and $\{0\}$ is a null set. So for $\operatorname{Re}s>-1$,
\begin{equation}\label{eq:Zpabs}
\int_{\Z_p}|x|_p^s\dd x=\sum_{k\ge0}p^{-ks}\,p^{-k}\Bigl(1-\frac1p\Bigr)=\frac{1-p^{-1}}{1-p^{-1-s}} .
\end{equation}

On the multiplicative group $\Q_p^\times$ we use
\[
\dd^\times x=\Bigl(1-\frac1p\Bigr)^{-1}\frac{\dd x}{|x|_p},
\]
normalised so that $\Z_p^\times$ has mass $1$. It is invariant under multiplication: for $a\in\Q_p^\times$, by
Proposition~\ref{prop:haarQp} applied to $g(x)=f(x)/|x|_p$,
\[
\int\frac{f(ax)}{|x|_p}\dd x=|a|_p\int\frac{f(ax)}{|ax|_p}\dd x=\int\frac{f(y)}{|y|_p}\dd y .
\]
On $\R^\times$ we use $\dd^\times x=\frac{\dd x}{|x|}$, which is invariant under multiplication for the same reason.

\begin{definition}
The \emph{local zeta integral} of a function $f$ on $\Q_v$ ($v$ a prime or $\infty$) is
$Z_v(f,s)=\int_{\Q_v^\times}f(x)|x|_v^s\dd^\times x$, when it converges.
\end{definition}

\begin{proposition}[the Euler factors are local integrals]\label{prop:localzeta}
For $\operatorname{Re}s>0$:
\[
Z_p(\mathbf 1_{\Z_p},s)=\frac1{1-p^{-s}},\qquad Z_\infty(e^{-\pi x^2},s)=\pi^{-s/2}\Gamma\Bigl(\frac s2\Bigr).
\]
\end{proposition}
\begin{proof}
By~\eqref{eq:Zpabs} with $s-1$ in place of $s$, $Z_p(\mathbf 1_{\Z_p},s)=(1-\frac1p)^{-1}\int_{\Z_p}|x|^{s-1}\dd x=(1-p^{-s})^{-1}$. At
infinity, by evenness and the substitution $u=\pi x^2$, $\frac{\dd x}x=\frac{\dd u}{2u}$:
\[
2\int_0^\infty e^{-\pi x^2}x^s\frac{\dd x}x=\int_0^\infty e^{-u}\Bigl(\frac u\pi\Bigr)^{s/2}\frac{\dd u}u=\pi^{-s/2}\Gamma\Bigl(\frac s2\Bigr).
\]
\end{proof}

Look at the product of all these local integrals: it is $\pi^{-s/2}\Gamma(\frac s2)\prod_p(1-p^{-s})^{-1}=\Lambda(s)$, Riemann's completed
zeta function, with the Gamma factor sitting at $\infty$ next to its brothers. The task of this chapter is to turn that observation
into a proof.

\subsection*{Integrals on \texorpdfstring{$\mathbb{A}_f$}{A\_f} and on \texorpdfstring{$\hat{\Z}^\times$}{Z-hat-times}}

The global theory only needs to integrate very simple functions, and we do it by hand.

\begin{definition}
$\mathcal S(\mathbb{A}_f)$ is the space of finite linear combinations of indicator functions $\mathbf 1_{b+N\hat{\Z}}$, with $b\in\mathbb{A}_f$ and $N\in\Q_{>0}$.
Its elements are the \emph{Schwartz--Bruhat functions} on $\mathbb{A}_f$.
\end{definition}

Every $f\in\mathcal S(\mathbb{A}_f)$ can be written $f=\sum_jc_j\mathbf 1_{b_j+M\hat{\Z}}$ with a single positive integer $M$ and pairwise disjoint
cosets. Indeed $N\hat{\Z}\supseteq a\hat{\Z}$ if $N=\frac ac$, so we may take $M$ to be a common multiple of the numerators of all the $N$'s;
then each $b+N\hat{\Z}$ is a finite union of cosets of $M\hat{\Z}$ (finite since $b+N\hat{\Z}$ is compact and the cosets are open and
disjoint). We set
\[
\int_{\mathbb{A}_f}f(x)\dd x=\sum_jc_j\,\frac1M .
\]
This does not depend on the representation: refining $M$ to a multiple $M'$ replaces each coset of $M\hat{\Z}$ by the $\frac{M'}M$ cosets of
$M'\hat{\Z}$ inside it ($\hat{\Z}/\frac{M'}M\hat{\Z}\cong\Z/\frac{M'}M\Z$ by density of $\Z$), each of mass $\frac1{M'}$, so the sum is unchanged; and two
representations have a common refinement, on which they agree since a function determines its values on each coset. The integral is
linear, positive, invariant under translation by any $a\in\mathbb{A}_f$ (translation permutes the cosets of $M\hat{\Z}$), and gives
$\int\mathbf 1_{b+N\hat{\Z}}=\frac1N$ for every rational $N>0$: if $N=\frac ac$, then $b+N\hat{\Z}$ is the union of the $c$ disjoint cosets
$b+\frac{aj}c+a\hat{\Z}$, $0\le j<c$, of mass $\frac1a$ each. In particular $\int\mathbf 1_{b+N\hat{\Z}}=\prod_p\mu_p(b_p+N\Z_p)$, since
$\mu_p(b_p+N\Z_p)=|N|_p$ and $\prod_p|N|_p=\frac1N$: this is the product of the local Haar measures. Multiplication by
$u\in\hat{\Z}^\times$ maps $b+N\hat{\Z}$ to $ub+N\hat{\Z}$, so $\int f(ux)\dd x=\int f(x)\dd x$.

On the compact group $\hat{\Z}^\times=\prod_p\Z_p^\times$ (the units of $\hat{\Z}$) we integrate locally constant functions. Such a function
$g$ factors through $(\Z/M\Z)^\times$ for some $M$: each $u$ has a neighbourhood $(u+M_u\hat{\Z})\cap\hat{\Z}^\times$ on which $g$ is constant;
finitely many of them cover $\hat{\Z}^\times$, and their $M_u$ divide some $M$; if $u\equiv u'$ modulo $M$, both lie in the same one of those
neighbourhoods, so $g(u)=g(u')$. We set
\[
\int_{\hat{\Z}^\times}g(u)\dd^\times u=\frac1{\varphi(M)}\sum_{a\in(\Z/M\Z)^\times}g(a).
\]
This does not depend on $M$: for $M\mid M'$, the reduction $(\Z/M'\Z)^\times\to(\Z/M\Z)^\times$ is a surjective homomorphism, so all its fibres
have $\frac{\varphi(M')}{\varphi(M)}$ elements. The integral is invariant under $u\mapsto vu$ ($v\in\hat{\Z}^\times$) and under $u\mapsto u^{-1}$,
which permute $(\Z/M\Z)^\times$, and $\int\dd^\times u=1$. For a Dirichlet character $\chi$ modulo $N$, viewed on $\hat{\Z}^\times$ by
$\tilde\chi(u)=\chi(u\bmod N)$, orthogonality gives $\int\tilde\chi\dd^\times u=1$ if $\chi$ is trivial and $0$ otherwise.

\begin{think}
Why not use the measure of $\mathbb{A}_f$ on $\hat{\Z}^\times$? Because it sees nothing: for every $y$, $\hat{\Z}^\times\subseteq\prod_{p\le y}\Z_p^\times\times\prod_{p>y}\Z_p$, a union of cosets of
$\bigl(\prod_{p\le y}p\bigr)\hat{\Z}$ of total mass $\prod_{p\le y}(1-\frac1p)$, which tends to $0$ (Tome~5: $\sum\frac1p$ diverges). So the units of
$\hat{\Z}$ form a null set, the adelic version of ``a random integer is coprime to every prime with probability $0$''. The multiplicative
measure has to be normalised separately, exactly as $\dd^\times x$ on $\Q_p$ carries the factor $(1-\frac1p)^{-1}$.
\end{think}

\section{Characters and Fourier analysis}

Fourier analysis on $\R$ uses the characters $x\mapsto e^{-2\pi ixy}$. On $\Q_p$ and $\mathbb{A}$ we need their analogues.

\begin{definition}
$\psi_p(x)=e^{2\pi i\{x\}_p}$ for $x\in\Q_p$, \quad $\psi_\infty(x)=e^{-2\pi ix}$ for $x\in\R$, \quad
$\psi_f(x)=\prod_p\psi_p(x_p)$ for $x\in\mathbb{A}_f$, \quad and $\psi(x)=\psi_\infty(x_\infty)\psi_f(x_f)$ for $x\in\mathbb{A}$.
\end{definition}

The product defining $\psi_f$ is finite, since $\psi_p(x_p)=1$ when $x_p\in\Z_p$. Each $\psi_p$ is a character of $(\Q_p,+)$: $\{x+y\}_p$
and $\{x\}_p+\{y\}_p$ both lie in $\Z[\frac1p]$ and differ by an element of $\Z_p$, hence by an element of $\Z_p\cap\Z[\frac1p]=\Z$. The sign
at infinity is the price of the next proposition.

\begin{proposition}\label{prop:psiQ}
\begin{enumerate}[label=(\alph*)]
\item $\psi$ is a continuous character of $\mathbb{A}$, trivial on $\Q$.
\item The kernel of $\psi_f$ is $\hat{\Z}$; more precisely $\psi_f(x)=1$ iff $x\in\hat{\Z}$. And $\psi_f(q)=e^{2\pi iq}$ for $q\in\Q$.
\end{enumerate}
\end{proposition}
\begin{proof}
(b) If $\psi_f(x)=1$ then $r=\sum_p\{x_p\}_p\in\Z$. Write $\{x_p\}_p=\frac{a_p}{p^{k_p}}$ with $p\nmid a_p$ when $k_p\ge1$. Multiplying
$r$ by $\prod_{\ell\ne p}\ell^{k_\ell}$ and reading the result in $\Z_p$, every term but the $p$-th is a $p$-adic integer, so
$\frac{a_p}{p^{k_p}}\in\Z_p$ and $k_p=0$: all fractional parts vanish, and $x\in\hat{\Z}$. The converse is clear. For $q\in\Q$,
Lemma~\ref{lem:AfQZ} gives $q-\sum_p\{q\}_p\in\Q\cap\hat{\Z}=\Z$, so $\psi_f(q)=e^{2\pi i\sum\{q\}_p}=e^{2\pi iq}$.

(a) $\psi$ is a product of characters, and $\psi(q)=e^{-2\pi iq}e^{2\pi iq}=1$. It is continuous: on the open set $\R\times(b+\hat{\Z})$ it
equals $x\mapsto e^{-2\pi ix_\infty}\psi_f(b)$, by (b).
\end{proof}

\subsection*{The Fourier transform on \texorpdfstring{$\mathbb{A}_f$}{A\_f}}

For $f\in\mathcal S(\mathbb{A}_f)$ and $y\in\mathbb{A}_f$, $x\mapsto f(x)\psi_f(xy)$ is again in $\mathcal S(\mathbb{A}_f)$ (on each coset of a small enough $M\hat{\Z}$
both factors are constant, by Proposition~\ref{prop:psiQ}(b)), and we define
\[
\hat f(y)=\int_{\mathbb{A}_f}f(x)\psi_f(xy)\dd x .
\]

\begin{proposition}\label{prop:FTAf}
For $b\in\mathbb{A}_f$ and $N\in\Q_{>0}$,
\[
\widehat{\mathbf 1_{b+N\hat{\Z}}}(y)=\psi_f(by)\,\frac1N\,\mathbf 1_{N^{-1}\hat{\Z}}(y).
\]
Hence $\hat f\in\mathcal S(\mathbb{A}_f)$ and $\hat{\hat f}(x)=f(-x)$ for every $f\in\mathcal S(\mathbb{A}_f)$.
\end{proposition}
\begin{proof}
Substituting $x=b+x'$ (translation invariance), $\widehat{\mathbf 1_{b+N\hat{\Z}}}(y)=\psi_f(by)I$ with $I=\int_{N\hat{\Z}}\psi_f(x'y)\dd x'$. If
$Ny\in\hat{\Z}$, the integrand is $1$ on $N\hat{\Z}$ (since $x'y\in Ny\hat{\Z}\subseteq\hat{\Z}$), and $I=\frac1N$. If $Ny\notin\hat{\Z}$, then
$h=N$ satisfies $\psi_f(hy)\ne1$ (Proposition~\ref{prop:psiQ}(b)), and translating by $h\in N\hat{\Z}$ gives $I=\psi_f(hy)I$, so $I=0$.

The function $y\mapsto\psi_f(by)\mathbf 1_{N^{-1}\hat{\Z}}(y)$ is in $\mathcal S(\mathbb{A}_f)$: it vanishes outside the compact open set $N^{-1}\hat{\Z}$
and is constant on cosets of $c\hat{\Z}$ whenever $bc\hat{\Z}\subseteq\hat{\Z}$, which holds for a suitable positive integer $c$ since $b\hat{\Z}$ is
contained in some $c^{-1}\hat{\Z}$. Finally, by linearity it suffices to check inversion on $f=\mathbf 1_{b+N\hat{\Z}}$. Using the first part
with the roles of $b$ and $x$ exchanged,
\[
\hat{\hat f}(x)=\frac1N\int_{N^{-1}\hat{\Z}}\psi_f\bigl((b+x)y\bigr)\dd y=\frac1N\cdot N\cdot\mathbf 1_{N\hat{\Z}}(b+x)=f(-x).
\]
\end{proof}

At infinity we work with the space $\mathcal S_\infty$ of finite sums of functions $P(x)e^{-\pi\lambda x^2}$, with $P$ a complex polynomial
and $\lambda>0$. It is stable under $x\mapsto tx$ ($t>0$) and under the Fourier transform: the transform of $e^{-\pi\lambda x^2}$ is
$\lambda^{-1/2}e^{-\pi y^2/\lambda}$ (Tome~4, Theorem~8.3), and by Tome~4, Proposition~8.2(d), the transform of $xf(x)$ is $\frac i{2\pi}\hat f\,'$; so
the transform of $x^ke^{-\pi\lambda x^2}$ is $\bigl(\frac{i}{2\pi}\bigr)^k\frac{\dd^k}{\dd y^k}\bigl(\lambda^{-1/2}e^{-\pi y^2/\lambda}\bigr)$, again in
$\mathcal S_\infty$. For instance
\begin{equation}\label{eq:xgauss}
\widehat{xe^{-\pi x^2}}(y)=\frac{i}{2\pi}\cdot(-2\pi y)e^{-\pi y^2}=-i\,ye^{-\pi y^2}.
\end{equation}
(Any rapidly decreasing smooth functions would do; this class suffices for $\zeta$ and Dirichlet $L$-functions.)

\begin{definition}
$\mathcal S(\mathbb{A})$ is the space of finite sums of functions $f(x)=f_\infty(x_\infty)f_f(x_f)$ with $f_\infty\in\mathcal S_\infty$ and
$f_f\in\mathcal S(\mathbb{A}_f)$. The Fourier transform is defined on such products by $\hat f=\hat f_\infty\,\hat f_f$, and extended linearly.
\end{definition}

For a product function,
\[
\hat f(y)=\int_\R\int_{\mathbb{A}_f}f(x)\psi(xy)\dd x_f\dd x_\infty,
\]
so the definition does not depend on how $f$ is written as a
sum of products, and $\hat{\hat f}(x)=f(-x)$ on $\mathcal S(\mathbb{A})$ by Proposition~\ref{prop:FTAf} and Tome~4, Theorem~8.6.

\begin{theorem}[adelic Poisson summation]\label{thm:adelicpoisson}
For every $f\in\mathcal S(\mathbb{A})$,
\[
\sum_{q\in\Q}f(q)=\sum_{q\in\Q}\hat f(q),
\]
both series converging absolutely.
\end{theorem}
\begin{proof}
By linearity, take $f=f_\infty\otimes\mathbf 1_{b+N\hat{\Z}}$. We may assume $b\in\Q$: by Lemma~\ref{lem:AfQZ} and the density of $\Z$ in
$\hat{\Z}$, the open set $b+N\hat{\Z}$ contains a rational, which we may use as centre. A rational $q$ lies in $b+N\hat{\Z}$ iff
$q-b\in\Q\cap N\hat{\Z}=N\Z$. So, with $g(t)=f_\infty(b+Nt)$,
\[
\sum_{q\in\Q}f(q)=\sum_{m\in\Z}f_\infty(b+Nm)=\sum_{m\in\Z}g(m).
\]
On the other side, $\hat f(q)=\hat f_\infty(q)\psi_f(bq)\frac1N\mathbf 1_{N^{-1}\hat{\Z}}(q)$ by Proposition~\ref{prop:FTAf}; the rationals in
$N^{-1}\hat{\Z}$ are the $\frac kN$, $k\in\Z$, and $\psi_f(\frac{bk}N)=e^{2\pi ibk/N}$ by Proposition~\ref{prop:psiQ}(b). So
\[
\sum_{q\in\Q}\hat f(q)=\frac1N\sum_{k\in\Z}e^{2\pi ibk/N}\hat f_\infty\Bigl(\frac kN\Bigr).
\]
By Tome~4, Proposition~8.2(a),(b), $\hat g(\xi)=\frac1Ne^{2\pi ib\xi/N}\hat f_\infty(\frac\xi N)$. Both $g$ and $\hat g$ are continuous
and $O((1+|t|)^{-2})$, so Tome~4, Theorem~8.9 gives $\sum_mg(m)=\sum_k\hat g(k)$, which is the claim.
\end{proof}

We shall use it in a scaled form. For $t>0$ and $u\in\hat{\Z}^\times$, write $(t,u)$ for the ad\`ele with components $t$ at infinity and
$u$ at the finite places, and $q(t,u)=(qt,qu)$.

\begin{corollary}\label{cor:scaledpoisson}
For $f\in\mathcal S(\mathbb{A})$, $t>0$ and $u\in\hat{\Z}^\times$,
\[
\sum_{q\in\Q}f\bigl(q(t,u)\bigr)=\frac1t\sum_{q\in\Q}\hat f\bigl(q(t^{-1},u^{-1})\bigr).
\]
\end{corollary}
\begin{proof}
Apply Theorem~\ref{thm:adelicpoisson} to $g(x)=f\bigl(x(t,u)\bigr)$, which lies in $\mathcal S(\mathbb{A})$ (multiplication by $u$ permutes the
cosets of $M\hat{\Z}$, and $\mathcal S_\infty$ is stable under $x\mapsto tx$). Its transform is $\hat g(y)=\frac1t\hat f\bigl(y(t^{-1},u^{-1})\bigr)$:
at infinity by Tome~4, Proposition~8.2(b), and on $\mathbb{A}_f$ by the substitution $x'=ux$, which preserves the integral.
\end{proof}

\section{The id\`eles of \texorpdfstring{$\Q$}{Q}}

\begin{definition}
The group of \emph{id\`eles} is
\[
\mathbb{A}^\times=\bigl\{x\in\mathbb{A}:\ x_v\ne0\text{ for all }v,\ x_p\in\Z_p^\times\text{ for all but finitely many }p\bigr\},
\]
the group of units of the ring $\mathbb{A}$. Its topology is generated by the sets $U_\infty\times\prod_{p\in S}U_p\times\prod_{p\notin S}\Z_p^\times$, with
$S$ finite and $U_v\subseteq\Q_v^\times$ open. The \emph{content} of an id\`ele is $|x|=\prod_v|x_v|_v$ (a finite product), and
$\mathbb{A}^1=\{x:|x|=1\}$.
\end{definition}

This topology is finer than the one induced by $\mathbb{A}$, and it has to be: in the induced topology inversion is not continuous
(Exercise~\ref{ex:ideletop}). The content is a continuous homomorphism $\mathbb{A}^\times\to\R_{>0}$, and $|q|=1$ for $q\in\Q^\times$ by the product
formula, so $\Q^\times\subseteq\mathbb{A}^1$.

\begin{proposition}\label{prop:ideleQ}
Every id\`ele can be written uniquely as $x=q\cdot(t,u)$ with $q\in\Q^\times$, $t>0$, $u\in\hat{\Z}^\times$, and then $|x|=t$. Hence
\[
\mathbb{A}^\times/\Q^\times\cong\R_{>0}\times\hat{\Z}^\times,\qquad \mathbb{A}^1/\Q^\times\cong\hat{\Z}^\times .
\]
In particular $\mathbb{A}^1/\Q^\times$ is compact.
\end{proposition}
\begin{proof}
Let $q=\operatorname{sign}(x_\infty)\prod_pp^{v_p(x_p)}$, a finite product. Then $q^{-1}x$ has positive real component and unit components
at every prime. If $q(t,u)=q'(t',u')$, then $q/q'$ is a positive rational that is a unit at every prime, so $q=q'$. Finally
$|(t,u)|=t\prod_p|u_p|_p=t$, and $|q|=1$.
\end{proof}

So the id\`ele class group of $\Q$ is the product of a line and a profinite group. In Chapter~10 we shall see that
$\hat{\Z}^\times=\lim_{\leftarrow}(\Z/N\Z)^\times$ is the Galois group of the maximal abelian extension of $\Q$: the id\`ele class group, built from
$\Q$ alone, knows all abelian extensions of $\Q$. That is the starting point of class field theory.

\section{Tate's thesis for \texorpdfstring{$\Q$}{Q}}

\begin{definition}
Let $f\in\mathcal S(\mathbb{A})$ and let $\chi$ be a Dirichlet character, viewed on $\hat{\Z}^\times$ as $\tilde\chi$. Put
\[
F_f(t,u)=\sum_{q\in\Q^\times}f\bigl(q(t,u)\bigr)\qquad(t>0,\ u\in\hat{\Z}^\times),
\]
and define the \emph{global zeta integral}
\[
Z(f,\chi,s)=\int_0^\infty\int_{\hat{\Z}^\times}F_f(t,u)\tilde\chi(u)\dd^\times u\ t^s\frac{\dd t}t .
\]
\end{definition}

By Proposition~\ref{prop:ideleQ}, $(t,u)$ runs over a fundamental domain for $\mathbb{A}^\times/\Q^\times$, and summing over $q$ unfolds the integral:
$Z(f,\chi,s)$ is the integral of $f(x)\tilde\chi(x)|x|^s$ over all of $\mathbb{A}^\times$, for the product of the measures $\dd^\times x_v$. We do not
need this interpretation, and we shall check it on examples instead.

\begin{lemma}\label{lem:Fbounds}
Let $f\in\mathcal S(\mathbb{A})$. There is a positive integer $L$ such that $F_f(t,u)$ depends on $u$ only through $u\bmod L$. For every $A>0$
there is $C$ with $|F_f(t,u)|\le Ct^{-A}$ for $t\ge1$, and $|F_f(t,u)|\le Ct^{-1}$ for $0<t\le1$.
\end{lemma}
\begin{proof}
Take $f=f_\infty\otimes f_f$ (sums are handled term by term). The support of $f_f$ is compact, so it is contained in $M^{-1}\hat{\Z}$ for some
positive integer $M$ (the $M^{-1}\hat{\Z}$ are open and increasing, and cover $\mathbb{A}_f$), and $f_f$ is constant on the cosets of some $L_0\hat{\Z}$, $L_0$ a
positive integer. Since $u$ is a unit, $qu\in M^{-1}\hat{\Z}$ forces $q\in M^{-1}\hat{\Z}\cap\Q=M^{-1}\Z$, so
\[
F_f(t,u)=\sum_{k\ne0}f_\infty\Bigl(\frac{kt}M\Bigr)f_f\Bigl(\frac{ku}M\Bigr),
\]
and $u\equiv u'\pmod{L_0M}$ implies $\frac{ku}M-\frac{ku'}M\in L_0\hat{\Z}$, so $L=L_0M$ works. Now $|f_f|\le c$ and
$|f_\infty(x)|\le c_A(1+|x|)^{-A-2}$ for every $A$. For $t\ge1$, $1+\frac{|k|t}M$ is at least $\frac{|k|t}M$ and at least $1+\frac{|k|}M$, so
$\sum_{k\ne0}(1+\frac{|k|t}M)^{-A-2}\le M^At^{-A}\sum_{k\ne0}(1+\frac{|k|}M)^{-2}=O(t^{-A})$. For $t\le1$, compare with an integral: $\sum_{k\ne0}(1+|k|t/M)^{-2}\le2\int_0^\infty(1+xt/M)^{-2}\dd x=\frac{2M}t$.
\end{proof}

\begin{theorem}[Tate, 1950, for $\Q$]\label{thm:tate}
Let $f\in\mathcal S(\mathbb{A})$ and $\chi$ a Dirichlet character. The integral $Z(f,\chi,s)$ converges absolutely for $\operatorname{Re}s>1$ and extends
to a meromorphic function on $\C$ satisfying
\[
Z(f,\chi,s)=Z(\hat f,\overline\chi,1-s).
\]
If $\chi$ is non-trivial it is entire. If $\chi$ is trivial, its only poles are simple poles at $s=1$, with residue $\hat f(0)$, and at
$s=0$, with residue $-f(0)$.
\end{theorem}
\begin{proof}
Let $\sigma=\operatorname{Re}s$. By Lemma~\ref{lem:Fbounds}, the inner integral $\Phi(t)=\int F_f(t,u)\tilde\chi(u)\dd^\times u$ is a finite average,
continuous in $t$, with $|\Phi(t)|\le Ct^{-1}$ near $0$ and $\le Ct^{-A}$ near $\infty$. So $\int_0^\infty|\Phi(t)|t^{\sigma-1}\dd t$ converges for
$\sigma>1$, and $\int_1^\infty\Phi(t)t^{s-1}\dd t$ converges for all $s$ and is entire (Tome~4, Theorem~9.15). Write $\delta_\chi=1$ if $\chi$ is
trivial and $0$ otherwise.

Corollary~\ref{cor:scaledpoisson}, with the terms $q=0$ separated, reads
\[
F_f(t,u)+f(0)=\frac1t\Bigl(F_{\hat f}\bigl(t^{-1},u^{-1}\bigr)+\hat f(0)\Bigr).
\]
Multiply by $\tilde\chi(u)t^{s-1}$ and integrate over $\hat{\Z}^\times\times\left]0,1\right]$, with $\sigma>1$. The first term gives the part
$\int_0^1\Phi(t)t^{s-1}\dd t$ of $Z(f,\chi,s)$. Moving $f(0)$ to the right, the constant terms give
$\delta_\chi\bigl(\hat f(0)\int_0^1t^{s-2}\dd t-f(0)\int_0^1t^{s-1}\dd t\bigr)=\delta_\chi\bigl(\frac{\hat f(0)}{s-1}-\frac{f(0)}s\bigr)$. In the remaining
term substitute $u\mapsto u^{-1}$ (which preserves $\dd^\times u$ and turns $\tilde\chi$ into $\tilde{\overline\chi}$) and $t\mapsto t^{-1}$:
\[
\int_0^1\!\!\int F_{\hat f}(t^{-1},u^{-1})\tilde\chi(u)\dd^\times u\,t^{s-2}\dd t=\int_1^\infty\!\!\int F_{\hat f}(t,u)\tilde{\overline\chi}(u)\dd^\times u\,t^{1-s}\frac{\dd t}t .
\]
Altogether, for $\sigma>1$,
\begin{equation}\label{eq:tatesplit}
\begin{aligned}
Z(f,\chi,s)&=\int_1^\infty\!\!\int F_f(t,u)\tilde\chi(u)\dd^\times u\,t^s\frac{\dd t}t\\
&\quad+\int_1^\infty\!\!\int F_{\hat f}(t,u)\tilde{\overline\chi}(u)\dd^\times u\,t^{1-s}\frac{\dd t}t\\
&\quad+\delta_\chi\Bigl(\frac{\hat f(0)}{s-1}-\frac{f(0)}s\Bigr).
\end{aligned}
\end{equation}
The two integrals are entire, so the right-hand side is the announced meromorphic continuation, with the announced poles and residues.
Apply~\eqref{eq:tatesplit} to $(\hat f,\overline\chi,1-s)$: since $\hat{\hat f}(x)=f(-x)$ and $q\mapsto-q$ permutes $\Q^\times$, $F_{\hat{\hat f}}=F_f$ and
$\hat{\hat f}(0)=f(0)$; the two integrals are exchanged, and the constant becomes $\delta_\chi\bigl(\frac{f(0)}{-s}-\frac{\hat f(0)}{1-s}\bigr)$, the same as
before. So both sides agree.
\end{proof}

Notice the structure of the proof: the functional equation comes from Poisson summation, and the poles come from the single term
$q=0$ that the sum over $\Q^\times$ leaves out. That is all.

\begin{corollary}[Riemann]\label{cor:riemannFE}
$\Lambda(s)=\pi^{-s/2}\Gamma(\frac s2)\zeta(s)$ extends meromorphically to $\C$, with simple poles at $1$ and $0$ of residues $1$ and $-1$, and
$\Lambda(s)=\Lambda(1-s)$.
\end{corollary}
\begin{proof}
Take $f=e^{-\pi x_\infty^2}\otimes\mathbf 1_{\hat{\Z}}$ and $\chi$ trivial. Then $\hat f=f$ (Tome~4, Theorem~8.3, and Proposition~\ref{prop:FTAf}
with $b=0$, $N=1$), and $f(0)=\hat f(0)=1$. For $u\in\hat{\Z}^\times$, $qu\in\hat{\Z}$ iff $q\in\Z$, so $F_f(t,u)=\sum_{n\ne0}e^{-\pi n^2t^2}$, and for
$\sigma>1$, using Proposition~\ref{prop:localzeta},
\[
Z(f,1,s)=2\sum_{n\ge1}\int_0^\infty e^{-\pi n^2t^2}t^s\frac{\dd t}t=\sum_{n\ge1}n^{-s}\,\pi^{-s/2}\Gamma\Bigl(\frac s2\Bigr)=\Lambda(s).
\]
Theorem~\ref{thm:tate} finishes the proof.
\end{proof}

The global integral is the product of the local ones of Proposition~\ref{prop:localzeta}, as promised: the Euler product is the factorisation
of an adelic integral over a product space. Exercise~\ref{ex:theta} shows that, unwound, the proof is exactly Riemann's.

Now the Dirichlet $L$-functions. Let $\chi$ be a primitive character modulo $N>1$, and let $a\in\{0,1\}$ with $\chi(-1)=(-1)^a$ (so $a=0$ for even
$\chi$ and $a=1$ for odd $\chi$). Put
\[
\Lambda(s,\chi)=\pi^{-(s+a)/2}\,\Gamma\Bigl(\frac{s+a}2\Bigr)L(s,\chi),
\]
and let $\phi_\chi\in\mathcal S(\mathbb{A}_f)$ be the function equal to $\overline\chi(x\bmod N)$ on $\hat{\Z}$ (with $\overline\chi=0$ on non-units modulo $N$) and to
$0$ outside $\hat{\Z}$: it is $\sum_{c\bmod N}\overline\chi(c)\mathbf 1_{c+N\hat{\Z}}$.

\begin{lemma}\label{lem:phichi}
$\hat{\phi}_\chi(y)=\frac{\tau(\overline\chi)}N\,\phi_{\overline\chi}(Ny)$, where $\phi_{\overline\chi}(z)=\chi(z\bmod N)\mathbf 1_{\hat{\Z}}(z)$.
\end{lemma}
\begin{proof}
By Proposition~\ref{prop:FTAf},
\[
\hat{\phi}_\chi(y)=\frac1N\sum_{c\bmod N}\overline\chi(c)\psi_f(cy)\mathbf 1_{N^{-1}\hat{\Z}}(y).
\]
For $y\in N^{-1}\hat{\Z}$ write $Ny=m+Nw$ with
$m\in\Z$, $w\in\hat{\Z}$ (density of $\Z$); then $cy=\frac{cm}N+cw$ and $\psi_f(cy)=e^{2\pi icm/N}$ by Proposition~\ref{prop:psiQ}(b). By Tome~5,
Theorem~3.5 applied to the primitive character $\overline\chi$, $\sum_c\overline\chi(c)e^{2\pi icm/N}=\chi(m)\tau(\overline\chi)$, and $\chi(m)=\chi(Ny\bmod N)$.
\end{proof}

\begin{corollary}[functional equation of Dirichlet $L$-functions]\label{cor:dirichletFE}
For $\chi$ primitive modulo $N>1$, $\Lambda(s,\chi)$ extends to an entire function, and
\[
\Lambda(s,\overline\chi)=(-i)^a\,\tau(\overline\chi)\,N^{-s}\,\Lambda(1-s,\chi).
\]
Equivalently, $\Lambda_N(s,\chi)=N^{s/2}\Lambda(s,\chi)$ satisfies $\Lambda_N(s,\chi)=\varepsilon(\chi)\Lambda_N(1-s,\overline\chi)$ with
$\varepsilon(\chi)=\dfrac{\tau(\chi)}{i^a\sqrt N}$, a complex number of modulus $1$.
\end{corollary}
\begin{proof}
Let $f_\infty(x)=x^ae^{-\pi x^2}$, so that $\hat f_\infty=(-i)^af_\infty$ by Tome~4, Theorem~8.3 and~\eqref{eq:xgauss}, and $f=f_\infty\otimes\phi_\chi$.
Then $\hat f(x)=c\,f_\infty(x_\infty)\phi_{\overline\chi}(Nx_f)$ with $c=(-i)^a\tau(\overline\chi)/N$, by Lemma~\ref{lem:phichi}.

\emph{The left side.} $\phi_\chi(qu)\ne0$ forces $q=n\in\Z$, and then $\phi_\chi(nu)=\overline\chi(n)\overline{\tilde\chi(u)}$. So
$\int F_f(t,u)\tilde\chi(u)\dd^\times u=\sum_{n\ne0}\overline\chi(n)f_\infty(nt)$, since $|\tilde\chi(u)|=1$. The terms $n$ and $-n$ are equal, because
$\overline\chi(-n)f_\infty(-nt)=(-1)^a(-1)^a\,\overline\chi(n)f_\infty(nt)$. Hence for $\sigma>1$
\[
Z(f,\chi,s)=2\sum_{n\ge1}\overline\chi(n)\int_0^\infty(nt)^ae^{-\pi n^2t^2}t^s\frac{\dd t}t=L(s,\overline\chi)\cdot2\int_0^\infty e^{-\pi t^2}t^{s+a}\frac{\dd t}t,
\]
and the last integral is $\pi^{-(s+a)/2}\Gamma(\frac{s+a}2)$ as in Proposition~\ref{prop:localzeta}. So $Z(f,\chi,s)=\Lambda(s,\overline\chi)$.

\emph{The right side.} $\phi_{\overline\chi}(Nqu)\ne0$ forces $q=\frac mN$, $m\in\Z$, and then $\phi_{\overline\chi}(mu)=\chi(m)\tilde\chi(u)$. Integrating against
$\tilde{\overline\chi}(u)$ leaves $c\sum_{m\ne0}\chi(m)f_\infty(\frac{mt}N)$, and the substitution $t\mapsto Nt$ gives, exactly as above,
$Z(\hat f,\overline\chi,1-s)=c\,N^{1-s}\Lambda(1-s,\chi)$.

Theorem~\ref{thm:tate} ($\chi$ is non-trivial) gives the entire continuation and the first identity. Exchanging the roles of $\chi$
and $\overline\chi$ (both are primitive with the same $a$) and multiplying by $N^{s/2}$,
\[
\Lambda_N(s,\chi)=(-i)^a\tau(\chi)N^{-1/2}\Lambda_N(1-s,\overline\chi),
\]
and $(-i)^a=i^{-a}$. Finally $|\tau(\chi)|=\sqrt N$ (Tome~5, Theorem~3.5).
\end{proof}

The \emph{root number} $\varepsilon(\chi)$ is a deep invariant. For a real character it equals $1$ (Exercise~\ref{ex:rootnumber}): this is Gauss's
determination of the sign of the quadratic Gauss sum, which took him four years. For complex characters it is a mysterious point on the
unit circle; the notebook computes a few and checks them against PARI.

\begin{think}
In Riemann's formula the factor $\pi^{-s/2}\Gamma(\frac s2)$ looks like something added by hand to make the equation symmetric. In Tate's proof
it is the local integral at $\infty$, computed with the Gaussian, exactly as $(1-p^{-s})^{-1}$ is the local integral at $p$ computed with
$\mathbf 1_{\Z_p}$. Why the Gaussian and $\mathbf 1_{\Z_p}$? Because each is its own Fourier transform: they are the ``unramified'' functions, and
the root number of $\zeta$ is $1$ because they are self-dual. For $L(s,\chi)$ the function at the primes dividing $N$ is not self-dual, and the
Gauss sum measures by how much.
\end{think}

\section{Number fields: compactness and its shadows}

Let $K$ be a number field of degree $n$ and signature $(r_1,r_2)$. For each place $v$ of $K$ let $K_v$ be the completion (Chapter~6), with its
normalised absolute value: $|x|_P=N(P)^{-v_P(x)}$ at a prime $P$, $|x|_v=|\sigma(x)|$ at a real place, $|x|_v=|\tau(x)|^2$ at a complex place;
$\mathcal O_v$ is the valuation ring at a finite place. The ad\`eles $\mathbb{A}_K$ and id\`eles $\mathbb{A}_K^\times$ are defined exactly as for $\Q$: families
$(x_v)$ with $x_v\in\mathcal O_v$ (respectively $x_v\in\mathcal O_v^\times$) for all but finitely many $v$, with the same restricted-product topologies.
The \emph{content} is $|x|=\prod_v|x_v|_v$, and $\mathbb{A}_K^1=\{|x|=1\}$ contains $K^\times$ by the product formula
(Proposition~\ref{prop:productformula}). To an id\`ele $x$ we attach the fractional ideal $\operatorname{id}(x)=\prod_PP^{v_P(x_P)}$; then
$N(\operatorname{id}(x))=\prod_P|x_P|_P^{-1}$.

\begin{lemma}[adelic Minkowski]\label{lem:adelicmink}
Let $C_K=(\frac2\pi)^{r_2}\sqrt{|d_K|}$. For every id\`ele $x$ with $|x|\ge C_K$ there is $\alpha\in K^\times$ with $|\alpha|_v\le|x_v|_v$ for every place $v$.
\end{lemma}
\begin{proof}
At the finite places the condition says $v_P(\alpha)\ge v_P(x_P)$ for all $P$, that is $\alpha\in I=\operatorname{id}(x)$. At the infinite places it says
$|\sigma(\alpha)|\le|x_v|$ at real places and $|\tau(\alpha)|\le|x_v|^{1/2}$ at complex places, where $|x_v|$ is the usual absolute value. The set $S$ of
points of $\R^{r_1}\times\C^{r_2}$ satisfying these bounds is compact, convex and symmetric, of volume
\[
\prod_{\text{real}}2|x_v|\prod_{\text{complex}}\pi|x_v|=2^{r_1}\pi^{r_2}\prod_{v\mid\infty}|x_v|_v .
\]
The lattice $j(I)$ has covolume $2^{-r_2}\sqrt{|d_K|}N(I)$ (Proposition~\ref{prop:covol}, extended to fractional ideals by scaling: if $cI$ is an
integral ideal, $c\in\Z_{>0}$, both sides are multiplied by $c^n$). Minkowski's theorem (Theorem~\ref{thm:minkowski}(b)) gives a non-zero $\alpha\in I$
with $j(\alpha)\in S$ as soon as
\[
2^{r_1}\pi^{r_2}\prod_{v\mid\infty}|x_v|_v\ge2^n2^{-r_2}\sqrt{|d_K|}\prod_P|x_P|_P^{-1},
\]
that is $|x|\ge2^{r_1+r_2}\pi^{-r_2}2^{-r_1}\sqrt{|d_K|}=C_K$.
\end{proof}

\begin{theorem}[Fujisaki]\label{thm:fujisaki}
$\mathbb{A}_K^1/K^\times$ is compact.
\end{theorem}
\begin{proof}
Choose an id\`ele $\theta$ with $\theta_v=1$ at every finite place and $|\theta|>C_K$ (make one archimedean component large). Let
\[
B=\bigl\{y\in\mathbb{A}_K^1:\ |y_v|_v\le|\theta_v|_v\text{ for every }v\bigr\}.
\]
\emph{$B$ meets every class.} Let $x\in\mathbb{A}_K^1$. Then $|x^{-1}\theta|=|\theta|>C_K$, so Lemma~\ref{lem:adelicmink} gives $\alpha\in K^\times$ with
$|\alpha|_v\le|x_v^{-1}\theta_v|_v$, that is $|\alpha x_v|_v\le|\theta_v|_v$, for all $v$; and $\alpha x\in\mathbb{A}_K^1$. So $\alpha x\in B$.

\emph{$B$ is compact.} For $y\in B$ and every $v$, since $\prod_w|y_w|_w=1$,
\[
|y_v|_v=\prod_{w\ne v}|y_w|_w^{-1}\ge\prod_{w\ne v}|\theta_w|_w^{-1}=\frac{|\theta_v|_v}{|\theta|}.
\]
At a finite place $P$ this says $|\theta|^{-1}\le|y_P|_P\le1$; if $N(P)>|\theta|$ the only value of $|\cdot|_P$ in this range is $1$, so $y_P\in\mathcal O_P^\times$. Let
$S$ be the finite set of archimedean places and of the $P$ with $N(P)\le|\theta|$. Then $B$ is contained in
\[
E=\prod_{v\in S}\Bigl\{z\in K_v:\ \frac{|\theta_v|_v}{|\theta|}\le|z|_v\le|\theta_v|_v\Bigr\}\times\prod_{v\notin S}\mathcal O_v^\times ,
\]
a product of compact subsets of the $K_v^\times$, which is compact in the id\`ele topology (Tychonoff), and $B$ is closed in $E$ (the conditions are
closed, and the content is continuous). So $B$ is compact, and its image $\mathbb{A}_K^1/K^\times$ is compact.
\end{proof}

\begin{corollary}\label{cor:hfinite}
The class group of $K$ is finite.
\end{corollary}
\begin{proof}
The map $x\mapsto\operatorname{id}(x)$ from $\mathbb{A}_K^1$ to the group of fractional ideals is a surjective homomorphism: any ideal is $\operatorname{id}(x)$ for an
id\`ele with the right finite components, and an archimedean component can be rescaled to make the content $1$. It sends $K^\times$ onto the principal
ideals, so it induces a surjection $\mathbb{A}_K^1/K^\times\to\operatorname{Cl}_K$. Its kernel contains the image of $H=\mathbb{A}_K^1\cap\bigl(\prod_{v\mid\infty}K_v^\times\times
\prod_P\mathcal O_P^\times\bigr)$, an open subgroup. The cosets of the image of $H$ form a cover of the compact group $\mathbb{A}_K^1/K^\times$ by disjoint open sets, so they
are finitely many, and $\operatorname{Cl}_K$ is finite.
\end{proof}

\begin{corollary}[Dirichlet's unit theorem, again]\label{cor:unitsagain}
$\ell(\mathcal O_K^\times)$ is a lattice in the hyperplane $\{\sum_vy_v=0\}$ of $\R^{r_1+r_2}$; so $\mathcal O_K^\times$ has rank $r_1+r_2-1$.
\end{corollary}
\begin{proof}
Here $\ell(x)=(\log|x|_v)_{v\mid\infty}$, and $\ell(\mathcal O_K^\times)$ is a discrete subgroup of the hyperplane $W$, by the product formula and
Proposition~\ref{prop:logkernel}. Keep $\theta$ and $B$ from the proof of Theorem~\ref{thm:fujisaki}. Let $y=(y_v)_{v\mid\infty}$ with $\prod_v|y_v|_v=1$, and
let $x$ be the id\`ele with these archimedean components and $1$ at every finite place; $x\in\mathbb{A}_K^1$. There is $\alpha\in K^\times$ with $\alpha x\in B$.
At the finite places this says $|\alpha|_P\le1$, so $\alpha\in\mathcal O_K$, and $|\alpha|_P\ge|\theta|^{-1}$, so $N(\alpha\mathcal O_K)\le|\theta|$. There are
finitely many integral ideals of norm at most $|\theta|$; let $\alpha_1,\dots,\alpha_m$ generate the principal ones. Then $\alpha=\varepsilon\alpha_i$ with
$\varepsilon\in\mathcal O_K^\times$ and some $i$. At the infinite places, $\frac{|\theta_v|_v}{|\theta|}\le|\alpha y_v|_v\le|\theta_v|_v$. Taking logarithms,
\[
\ell(y)=-\ell(\varepsilon)-\ell(\alpha_i)+b,\qquad b\in\prod_{v\mid\infty}\bigl[\log|\theta_v|_v-\log|\theta|,\ \log|\theta_v|_v\bigr].
\]
As $\ell(y)$ runs over all of $W$, every point of $W$ lies within bounded distance of $\ell(\mathcal O_K^\times)$. A discrete subgroup is generated by
linearly independent vectors (Lemma~\ref{lem:discrete}); if they did not span $W$, a vector $w\in W$ orthogonal to their span would give points $\lambda w$,
$\lambda\to\infty$, arbitrarily far from $\ell(\mathcal O_K^\times)$. So $\ell(\mathcal O_K^\times)$ is a lattice in $W$, of rank $\dim W=r_1+r_2-1$, and the kernel
$\mu_K$ of $\ell$ is finite (Proposition~\ref{prop:logkernel}).
\end{proof}

So the two finiteness theorems of Chapters~3 and~4 are consequences of one theorem, the compactness of $\mathbb{A}_K^1/K^\times$. In fact they are
equivalent to it.

\begin{proposition}\label{prop:fujisakiconverse}
The finiteness of $h_K$ and Dirichlet's unit theorem imply that $\mathbb{A}_K^1/K^\times$ is compact.
\end{proposition}
\begin{proof}
Let $H=\mathbb{A}_K^1\cap\bigl(\prod_{v\mid\infty}K_v^\times\times\prod_P\mathcal O_P^\times\bigr)=K_\infty^1\times\prod_P\mathcal O_P^\times$, where $K_\infty^1$ is the set of
archimedean families with $\prod_{v\mid\infty}|y_v|_v=1$. The kernel of $\mathbb{A}_K^1\to\operatorname{Cl}_K$ is $K^\times H$: if $\operatorname{id}(x)=\alpha\mathcal O_K$, then
$\alpha^{-1}x$ has unit components at every finite place and content $1$. So $\mathbb{A}_K^1=\bigcup_{i=1}^{h}x_iK^\times H$ for representatives $x_i$. By the unit
theorem $W=\ell(\mathcal O_K^\times)+P$ for a compact parallelotope $P$, so $K_\infty^1=\mathcal O_K^\times\cdot\ell^{-1}(P)$, and $\ell^{-1}(P)\cap K_\infty^1$ is
compact (the absolute values are bounded above and below; at complex places the argument runs over a circle). Hence $H=\mathcal O_K^\times C$ with $C$ compact,
$\mathbb{A}_K^1=\bigcup_ix_iK^\times C$, and $\mathbb{A}_K^1/K^\times$ is the image of the compact set $\bigcup_ix_iC$.
\end{proof}

\begin{remark}[Hecke's functional equation]
With compactness in hand, Tate's proof goes through for $K$ word for word: Haar measures on each $K_v$, a character $\psi_K=\psi\circ\operatorname{Tr}_{K/\Q}$
for which $K$ is discrete in $\mathbb{A}_K$ with compact quotient, adelic Poisson summation, and the splitting of the zeta integral at content $1$. The result is
Hecke's theorem: with $\Gamma_\R(s)=\pi^{-s/2}\Gamma(\frac s2)$ and $\Gamma_\C(s)=2(2\pi)^{-s}\Gamma(s)$, the function
\[
\Lambda_K(s)=|d_K|^{s/2}\,\Gamma_\R(s)^{r_1}\Gamma_\C(s)^{r_2}\,\zeta_K(s)
\]
extends meromorphically to $\C$, with simple poles at $0$ and $1$ only, and $\Lambda_K(s)=\Lambda_K(1-s)$. The residue at $s=1$ is $\frac{2^{r_1}2^{r_2}h_KR_K}{w_K}$,
which is Theorem~\ref{thm:cnf} multiplied by $|d_K|^{1/2}\Gamma_\R(1)^{r_1}\Gamma_\C(1)^{r_2}=\sqrt{|d_K|}\,\pi^{-r_2}$. In Tate's proof this residue is, up to the normalisation of the measures, the volume of
$\mathbb{A}_K^1/K^\times$, so the class number formula becomes the computation of a volume. Exercise~\ref{ex:hecke} proves the quadratic
case from the corollaries above.
\end{remark}

\section{Exercises}

\begin{exercise}[$\star$]\label{ex:strongapprox}
\begin{enumerate}[label=(\alph*)]
\item Show that $\Q$ is dense in $\mathbb{A}_f$.
\item Deduce: given distinct primes $p_1,\dots,p_r$, elements $x_i\in\Q_{p_i}$ and $k\ge1$, there is $q\in\Q$ with $|q-x_i|_{p_i}\le p_i^{-k}$ for all $i$ and
$q\in\Z_p$ for every other prime $p$.
\item Show that $\Q$ is not dense in $\mathbb{A}$, but that $\Q+(\R\times\{0\})$ is.
\end{enumerate}
\end{exercise}

\begin{exercise}[$\star$]\label{ex:ftzp}
On $\Q_p$ define $\hat f(y)=\int f(x)\psi_p(xy)\dd x$ for $f$ a linear combination of indicators of balls.
\begin{enumerate}[label=(\alph*)]
\item Show that $\widehat{\mathbf 1_{b+p^k\Z_p}}(y)=\psi_p(by)p^{-k}\mathbf 1_{p^{-k}\Z_p}(y)$; in particular $\widehat{\mathbf 1_{\Z_p}}=\mathbf 1_{\Z_p}$.
\item Compute the transform of $\mathbf 1_{\Z_p^\times}$, and check that transforming twice gives $\mathbf 1_{\Z_p^\times}$ back.
\end{enumerate}
\end{exercise}

\begin{exercise}[$\star$]\label{ex:theta}
Let $f=e^{-\pi x_\infty^2}\otimes\mathbf 1_{\hat{\Z}}$ and $\theta(y)=\sum_{n\in\Z}e^{-\pi n^2y}$. Show that $F_f(t,u)=\theta(t^2)-1$, that
Corollary~\ref{cor:scaledpoisson} becomes $\theta(1/y)=\sqrt y\,\theta(y)$, and that~\eqref{eq:tatesplit} becomes, after the substitution $y=t^2$, the
formula of Tome~4, Theorem~10.8.
\end{exercise}

\begin{exercise}[$\star$]\label{ex:rootnumber}
\begin{enumerate}[label=(\alph*)]
\item Compute $\tau(\chi)$ and $\varepsilon(\chi)$ for the real primitive characters modulo $3$, $4$, $5$ and $8$.
\item Let $p$ be an odd prime and $\chi=\bigl(\frac\cdot p\bigr)$. Show that $\tau(\chi)=\sum_{k=0}^{p-1}e^{2\pi ik^2/p}$, and deduce from Gauss's theorem
(Tome~4, Theorem~8.14) that $\varepsilon(\chi)=1$.
\end{enumerate}
\end{exercise}

\begin{exercise}[$\star\star$]\label{ex:hecke}
Let $p$ be an odd prime, $p^*=(-1)^{(p-1)/2}p$ and $K=\Q(\sqrt{p^*})$, so $d_K=p^*$ and $\zeta_K(s)=\zeta(s)L(s,\chi)$ with $\chi=\bigl(\frac\cdot p\bigr)$
(Theorem~\ref{thm:zetaab}). Show that $\Lambda_K(s)=\Lambda(s)\Lambda_p(s,\chi)$, with the notation of the Remark on Hecke's functional equation, and deduce
$\Lambda_K(s)=\Lambda_K(1-s)$. (For $p\equiv3\pmod 4$ you will need the duplication formula.)
\end{exercise}

\begin{exercise}[$\star\star$]\label{ex:ideletop}
\begin{enumerate}[label=(\alph*)]
\item Let $p_n$ be the $n$-th prime and $x^{(n)}$ the id\`ele with $p_n$ at the place $p_n$ and $1$ elsewhere. Show that $x^{(n)}\to1$ in $\mathbb{A}$, but that
$(x^{(n)})^{-1}$ does not converge to $1$ in $\mathbb{A}$. So inversion is not continuous for the topology induced by $\mathbb{A}$.
\item Show that the id\`ele topology is the topology induced by the injection $\mathbb{A}^\times\to\mathbb{A}\times\mathbb{A}$, $x\mapsto(x,x^{-1})$.
\end{enumerate}
\end{exercise}

\begin{exercise}[$\star\star\star$]\label{ex:selfdual}
Show that every continuous homomorphism $\chi:(\Q_p,+)\to\{z\in\C:|z|=1\}$ is $x\mapsto\psi_p(ax)$ for a unique $a\in\Q_p$. ($\Q_p$ is \emph{self-dual}. The
same holds for $\mathbb{A}$, whose characters trivial on $\Q$ are the $\psi(q\,\cdot)$, $q\in\Q$: this is the deeper reason why Poisson summation sums over $\Q$
on both sides.)
\end{exercise}

\begin{beauty}
Riemann needed a theta function and a clever substitution; Hecke needed theta functions in many variables and the units and class group of the field;
Tate needed only Poisson summation on a single locally compact ring, in which a number field sits as a lattice. The real numbers stop being special:
$\R$ is one completion among infinitely many, the Gamma function is an Euler factor, and the class number and the units are the two halves of one compact
group. The same philosophy, pushed to non-abelian groups, became the Langlands programme, in which automorphic representations of $\mathrm{GL}_n(\mathbb{A})$
play the role that Tate's characters of $\mathbb{A}^\times/\Q^\times$ play here.
\end{beauty}

\section{Hints}

Try each exercise for a while before reading its hint, and read the hint before the solution.

\begin{hint}{ex:strongapprox}
(a) Lemma~\ref{lem:AfQZ} and the density of $\Z$ in $\hat{\Z}$. (c) Theorem~\ref{thm:Qdiscrete}.
\end{hint}
\begin{hint}{ex:ftzp}
Follow the proof of Proposition~\ref{prop:FTAf}; the kernel of $\psi_p$ is $\Z_p$. Write $\mathbf 1_{\Z_p^\times}=\mathbf 1_{\Z_p}-\mathbf 1_{p\Z_p}$.
\end{hint}
\begin{hint}{ex:theta}
$qu\in\hat{\Z}$ iff $q\in\Z$. In Corollary~\ref{cor:scaledpoisson} the two sides contain all $q\in\Q$, including $0$.
\end{hint}
\begin{hint}{ex:rootnumber}
$\varepsilon(\chi)=\frac{\tau(\chi)}{i^a\sqrt N}$. For (b), count the solutions of $k^2\equiv c$ modulo $p$.
\end{hint}
\begin{hint}{ex:hecke}
For $p\equiv1\pmod4$, $K$ is real and $\chi$ is even; for $p\equiv3\pmod4$, $K$ is imaginary and $\chi$ is odd. Show that $\Gamma_\R(s)\Gamma_\R(s+1)=\Gamma_\C(s)$.
Use Exercise~\ref{ex:rootnumber}.
\end{hint}
\begin{hint}{ex:ideletop}
(a) A basic neighbourhood of $1$ in $\mathbb{A}$ only constrains finitely many places, and requires $\Z_p$ elsewhere. (b) $x_p\in\Z_p$ and $x_p^{-1}\in\Z_p$ means
$x_p\in\Z_p^\times$.
\end{hint}
\begin{hint}{ex:selfdual}
First show that $\chi$ is trivial on some $p^k\Z_p$ (the disc $|z-1|<1$ contains no non-trivial subgroup of the circle). Then study $\chi$ on the finite cyclic
groups $p^{-j}\Z_p/p^k\Z_p$, and pass to the limit in $\Z_p$.
\end{hint}

\section{Solutions}

\begin{solution}{ex:strongapprox}
\idea Every ad\`ele is a rational plus an element of $\hat{\Z}$, and integers approximate elements of $\hat{\Z}$; the place $\infty$ is what makes $\Q$ discrete.
\writeup (a) A basic neighbourhood of $x\in\mathbb{A}_f$ contains $x+N\hat{\Z}$ for a positive integer $N$. Write $x=r+z$ with $r\in\Q$, $z\in\hat{\Z}$
(Lemma~\ref{lem:AfQZ}) and choose $n\in\Z$ with $z-n\in N\hat{\Z}$ (Chinese remainder theorem). Then $q=r+n\in x+N\hat{\Z}$.

(b) Apply (a) to the finite ad\`ele $x$ with components $x_i$ at $p_i$ and $0$ elsewhere, and to $N=\prod p_i^k$: the $q$ found satisfies
$q-x_{p}\in N\Z_p$ for every $p$, which is the claim.

(c) $\Q$ is discrete (Theorem~\ref{thm:Qdiscrete}), hence closed: if $q_n\to x$ then $q_n-q_m\to0$, so $q_n$ is eventually constant. As $\Q\ne\mathbb{A}$, it is not
dense. For $x=(x_\infty,x_f)\in\mathbb{A}$ and a neighbourhood $U\times(x_f+N\hat{\Z})$, take $q$ as in (a) for $x_f$; then $q+(x_\infty-q,0)$ has components $x_\infty$
and $q\in x_f+N\hat{\Z}$, so it lies in the neighbourhood.
\end{solution}

\begin{solution}{ex:ftzp}
\idea A character integrates to the volume over a subgroup on which it is trivial, and to $0$ otherwise.
\writeup (a) By translation invariance, $\hat f(y)=\psi_p(by)\int_{p^k\Z_p}\psi_p(xy)\dd x$. If $p^ky\in\Z_p$ the integrand is $1$ and the integral is $p^{-k}$.
Otherwise $\{p^ky\}_p\ne0$, so $h=p^k$ satisfies $\psi_p(hy)\ne1$, and translation by $h$ shows that the integral $I$ satisfies $I=\psi_p(hy)I$, so $I=0$. For
$b=0$, $k=0$: $\widehat{\mathbf 1_{\Z_p}}=\mathbf 1_{\Z_p}$.

(b) $\widehat{\mathbf 1_{\Z_p^\times}}=\mathbf 1_{\Z_p}-\frac1p\mathbf 1_{p^{-1}\Z_p}$: it is $1-\frac1p$ on $\Z_p$, $-\frac1p$ on $p^{-1}\Z_p\setminus\Z_p$, and $0$ elsewhere.
Transforming again with (a) ($k=-1$ gives $\widehat{\mathbf 1_{p^{-1}\Z_p}}=p\mathbf 1_{p\Z_p}$): $\mathbf 1_{\Z_p}-\frac1p\cdot p\mathbf 1_{p\Z_p}=\mathbf 1_{\Z_p^\times}$, which is
$f(-x)=f(x)$, as Proposition~\ref{prop:FTAf} predicts locally.
\end{solution}

\begin{solution}{ex:theta}
\idea Restricted to this $f$, every adelic object is a classical one.
\writeup For $u\in\hat{\Z}^\times$, $f(q(t,u))=e^{-\pi q^2t^2}\mathbf 1_{\hat{\Z}}(qu)$ is non-zero iff $q\in\Z$, so $F_f(t,u)=\sum_{n\ne0}e^{-\pi n^2t^2}=\theta(t^2)-1$.
Since $\hat f=f$, Corollary~\ref{cor:scaledpoisson} reads $\theta(t^2)=t^{-1}\theta(t^{-2})$; with $y=t^{-2}$ this is $\theta(1/y)=\sqrt y\,\theta(y)$. In
\eqref{eq:tatesplit}, with $\chi$ trivial and $f(0)=\hat f(0)=1$, $\Lambda(s)=\int_1^\infty(\theta(t^2)-1)(t^s+t^{1-s})\frac{\dd t}t+\frac1{s-1}-\frac1s$. Put $y=t^2$,
$\frac{\dd t}t=\frac{\dd y}{2y}$, and $\omega(y)=\sum_{n\ge1}e^{-\pi n^2y}=\frac12(\theta(y)-1)$:
\[
\Lambda(s)=-\frac1s-\frac1{1-s}+\int_1^\infty\omega(y)\bigl(y^{s/2}+y^{(1-s)/2}\bigr)\frac{\dd y}y,
\]
which is Tome~4, Theorem~10.8, line for line.
\end{solution}

\begin{solution}{ex:rootnumber}
\idea Real characters have real or purely imaginary Gauss sums, and the sign is the whole difficulty.
\writeup (a) Write $e(x)=e^{2\pi ix}$. Modulo $3$ ($\chi$ odd, $a=1$): $\tau=e(\frac13)-e(\frac23)=i\sqrt3$, so $\varepsilon=\frac{i\sqrt3}{i\sqrt3}=1$. Modulo $4$
($a=1$): $\tau=i-(-i)=2i$, $\varepsilon=\frac{2i}{2i}=1$. Modulo $5$ ($\chi(\pm1)=1$, $\chi(\pm2)=-1$, $a=0$): $\tau=2\cos\frac{2\pi}5-2\cos\frac{4\pi}5=\sqrt5$, since
$\cos\frac{2\pi}5=\frac{\sqrt5-1}4$ and $\cos\frac{4\pi}5=-\frac{\sqrt5+1}4$; $\varepsilon=1$. Modulo $8$ ($\chi(\pm1)=1$, $\chi(\pm3)=-1$, $a=0$):
$\tau=2\cos\frac\pi4-2\cos\frac{3\pi}4=2\sqrt2=\sqrt8$, $\varepsilon=1$.

(b) The number of $k$ modulo $p$ with $k^2\equiv c$ is $1+\chi(c)$. So $\sum_ke(\frac{k^2}p)=\sum_c(1+\chi(c))e(\frac cp)=\sum_c\chi(c)e(\frac cp)=\tau(\chi)$, as
$\sum_ce(\frac cp)=0$. Gauss's theorem gives $\tau(\chi)=\frac{(1+i)(1+i^{-p})}2\sqrt p$, which is $\sqrt p$ if $p\equiv1\pmod4$ and $i\sqrt p$ if $p\equiv3\pmod4$. Since
$\chi(-1)=(-1)^{(p-1)/2}$, $a=0$ in the first case and $a=1$ in the second, and $\varepsilon(\chi)=\frac{\tau(\chi)}{i^a\sqrt p}=1$ in both.
\end{solution}

\begin{solution}{ex:hecke}
\idea The Gamma factors of $\zeta$ and $L(s,\chi)$ multiply to the Gamma factors of $K$, and the root numbers multiply to $1$.
\writeup If $p\equiv1\pmod4$, $K$ is real ($r_1=2$), $\chi$ is even, and
\[
\Lambda_K(s)=p^{s/2}\Gamma_\R(s)^2\zeta(s)L(s,\chi)=\Lambda(s)\Lambda_p(s,\chi).
\] If $p\equiv3\pmod4$, $K$ is imaginary
($r_2=1$), $\chi$ is odd, $|d_K|=p$, and the duplication formula $\Gamma(z)\Gamma(z+\frac12)=2^{1-2z}\sqrt\pi\,\Gamma(2z)$, valid for real $z>0$ and hence for
$\operatorname{Re}z>0$ by analytic continuation, gives with $z=\frac s2$
\[
\Gamma_\R(s)\Gamma_\R(s+1)=\pi^{-s-\frac12}\,2^{1-s}\sqrt\pi\,\Gamma(s)=2(2\pi)^{-s}\Gamma(s)=\Gamma_\C(s).
\]
Since $\Lambda_p(s,\chi)=p^{s/2}\Gamma_\R(s+1)L(s,\chi)$ for odd $\chi$,
\[
\Lambda_K(s)=p^{s/2}\Gamma_\R(s)\Gamma_\R(s+1)\zeta(s)L(s,\chi)=\Lambda(s)\Lambda_p(s,\chi).
\]
In both cases $\chi=\overline\chi$ and $\varepsilon(\chi)=1$ (Exercise~\ref{ex:rootnumber}), so Corollaries~\ref{cor:riemannFE} and~\ref{cor:dirichletFE} give
$\Lambda_K(1-s)=\Lambda(1-s)\Lambda_p(1-s,\chi)=\Lambda(s)\Lambda_p(s,\chi)=\Lambda_K(s)$.
\end{solution}

\begin{solution}{ex:ideletop}
\idea The ad\`ele topology only watches finitely many places closely; the id\`ele topology also watches that inverses stay integral.
\writeup (a) A basic neighbourhood of $1$ in $\mathbb{A}$ has the form $U_\infty\times\prod_{p\in S}U_p\times\prod_{p\notin S}\Z_p$. For $n$ large, $p_n\notin S$, and $x^{(n)}$
has component $1$ at all places of $S\cup\{\infty\}$ and $p_n\in\Z_{p_n}$ at $p_n$, so it lies in the neighbourhood. But $(x^{(n)})^{-1}$ has component
$\frac1{p_n}\notin\Z_{p_n}$, so it never lies in $\R\times\hat{\Z}$, a neighbourhood of $1$.

(b) Let $\iota(x)=(x,x^{-1})$. If $W=U_\infty\times\prod_{S}U_p\times\prod_{p\notin S}\Z_p^\times$ is a basic id\`ele open set, then $\iota(W)=\iota(\mathbb{A}^\times)\cap(V\times V')$
with $V=U_\infty\times\prod_SU_p\times\prod_{p\notin S}\Z_p$ and $V'$ the same with the inverse sets $U_v^{-1}$, because $x_p,x_p^{-1}\in\Z_p$ iff $x_p\in\Z_p^\times$; so
$W$ is open for the induced topology. Conversely, let $V,V'$ be basic ad\`ele open sets and $y\in\iota^{-1}(V\times V')$. Let $S$ contain the places where $V$ or $V'$ is not
$\Z_p$. Then $W=\prod_{v\in S\cup\{\infty\}}\{z\in\Q_v^\times:z\in V_v,\ z^{-1}\in V'_v\}\times\prod_{p\notin S}\Z_p^\times$ is a basic id\`ele open set (each factor is open in
$\Q_v^\times$, inversion being continuous there) containing $y$ and contained in $\iota^{-1}(V\times V')$. So the two topologies coincide.
\end{solution}

\begin{solution}{ex:selfdual}
\idea A continuous character is trivial on a small ball, so on each larger ball it is a character of a finite cyclic group; compatible choices converge in $\Z_p$.
\writeup \emph{Step 1.} If $z\ne1$ lies on the circle with $|z-1|<1$, some power $z^m$ has real part $\le0$, hence $|z^m-1|\ge1$: so the disc $D=\{|z-1|<1\}$ contains no
non-trivial subgroup. By continuity some $p^k\Z_p$ is mapped into $D$; its image is a subgroup, hence trivial.

\emph{Step 2.} Fix $j\ge0$. The group $p^{-j}\Z_p/p^k\Z_p$ is cyclic of order $p^{j+k}$, generated by $p^{-j}$, and $\chi$ induces a character of it, so
$\chi(p^{-j})=e^{2\pi im_j/p^{j+k}}$ for some integer $m_j$. For $c\in\Z_p$ the character $x\mapsto\psi_p(cp^{-k}x)$ is also trivial on $p^k\Z_p$, and takes the
value $e^{2\pi ic/p^{j+k}}$ (with $c$ read modulo $p^{j+k}$) at $p^{-j}$. Choosing $c_j\equiv m_j\pmod{p^{j+k}}$, the two characters agree on a generator, hence on
$p^{-j}\Z_p$.

\emph{Step 3.} $\psi_p(cp^{-k}\cdot)$ is trivial on $p^{-j}\Z_p$ iff $cp^{-k-j}\in\Z_p$, so $c_{j+1}\equiv c_j\pmod{p^{j+k}}$: the $c_j$ form a Cauchy sequence, with
limit $c\in\Z_p$, and $c\equiv c_j\pmod{p^{j+k}}$ for all $j$. Hence $\chi(x)=\psi_p(ax)$ with $a=cp^{-k}$ for $x$ in every $p^{-j}\Z_p$, that is for all $x$.
Uniqueness: if $\psi_p(ax)=1$ for all $x$, then $a\Q_p\subseteq\Z_p$, so $a=0$.
\end{solution}

\section*{Sources and further reading}

\begin{itemize}
\item J.~T.~Tate, \emph{Fourier analysis in number fields and Hecke's zeta-functions}, Ph.D.\ thesis, Princeton, 1950; reprinted in J.~W.~S.~Cassels and
A.~Fr\"ohlich (eds.), \emph{Algebraic Number Theory}, Academic Press, 1967, 305--347. Still the best read on the subject.
\item C.~Chevalley, ``G\'en\'eralisation de la th\'eorie du corps de classes pour les extensions infinies'', \emph{J.\ Math.\ Pures Appl.}\ 15 (1936),
359--371. The first appearance of id\`eles.
\item K.~Iwasawa, ``Letter to J.~Dieudonn\'e'' (1952), in \emph{Zeta Functions in Geometry}, Adv.\ Stud.\ Pure Math.\ 21 (1992), 445--450. The parallel discovery.
\item D.~Ramakrishnan and R.~J.~Valenza, \emph{Fourier Analysis on Number Fields}, Springer GTM~186, 1999. Haar measure, restricted products and Tate's thesis
in full generality.
\item S.~Lang, \emph{Algebraic Number Theory}, 2nd ed., Springer GTM~110, 1994, Chapters~VII and~XIV. Fujisaki's theorem and Tate's thesis.
\item J.~Neukirch, \emph{Algebraic Number Theory}, Springer, 1999, Chapters~VI and~VII. Id\`eles, and Hecke $L$-functions with their functional equation.
\item E.~Hecke, ``Eine neue Art von Zetafunktionen und ihre Beziehungen zur Verteilung der Primzahlen'', \emph{Math.\ Z.}\ 1 (1918), 357--376, and 6 (1920),
11--51.
\item S.~Gelbart and S.~D.~Miller, ``Riemann's zeta function and beyond'', \emph{Bull.\ AMS} 41 (2004), 59--112. From Riemann through Tate to Langlands.
\end{itemize}
